\documentclass[11pt]{amsart}
\renewcommand{\bfdefault}{m}
\usepackage[T1]{fontenc}
\usepackage{lmodern}
\usepackage[letterpaper,textwidth=6in,textheight=8.5in,centering]{geometry}
\usepackage{amsmath,amssymb,amsthm,mathtools}
\usepackage[hidelinks]{hyperref}
\usepackage[nameinlink,noabbrev]{cleveref}
\newtheorem{theorem}{Theorem}[section]
\newtheorem{proposition}[theorem]{Proposition}
\newtheorem{lemma}[theorem]{Lemma}
\newtheorem{corollary}[theorem]{Corollary}
\theoremstyle{remark}
\newtheorem{remark}[theorem]{Remark}
\crefname{theorem}{Theorem}{Theorems}
\crefname{proposition}{Proposition}{Propositions}
\crefname{lemma}{Lemma}{Lemmas}
\crefname{corollary}{Corollary}{Corollaries}
\crefname{remark}{Remark}{Remarks}
\newcommand{\EE}{\mathbb E}
\newcommand{\PP}{\mathbb P}
\newcommand{\ZZ}{\mathbb Z}
\newcommand{\RR}{\mathbb R}
\DeclareMathOperator{\Exp}{Exp}
\DeclareMathOperator{\Ber}{Bernoulli}
\DeclareMathOperator*{\essinf}{ess\,inf}
\numberwithin{equation}{section}
\title{A flux plateau for site-disordered finite-capacity exclusion}
\author{The author}
\address{}

\date{}
\subjclass[2020]{60K35, 60K37, 82C22}
\keywords{Exclusion process, site disorder, flux plateau, last-passage percolation}

\begin{document}
\begin{abstract}
We consider totally asymmetric exclusion on the integer lattice with capacity
$K$ at each site and independent, identically distributed site rates $\alpha_i$.
For every finite $K$ and every nonconstant rate law with positive essential
infimum $r$, the flux has a nondegenerate flat interval about density $K/2$.
Its height is $rM_K$, where $M_K$ is the maximum homogeneous rate-one flux,
and a guaranteed half-width is
$(1/r-\EE[1/\alpha_0])rM_K$.
The proof extends Sly's exponential-delay comparison to the three-step growth
representation of finite-capacity exclusion. Homogeneous reflection symmetry
provides the supporting plane, and long intervals of slow sites determine
the maximal current.
\end{abstract}
\maketitle

\section{Introduction}

Consider a particle system on $\ZZ$ in which each site holds at most $K$
particles, where $K$ is a fixed positive integer. A clock at site $i$ rings
at rate $\alpha_i$. At a ring, one particle moves from $i$ to $i+1$ if the
departure site is nonempty and the receiving site has fewer than $K$
particles. The rates form an independent, identically distributed
environment $\alpha=(\alpha_i)_{i\in\ZZ}$ with common law $Q$. We assume
that $Q$ is supported on $[r,1]$, where
\[
 r=\essinf\alpha_0>0.
\]
There is one clock per departure site; its rate does not increase with the
occupation of that site.

Sepp\"al\"ainen introduces this model in equation (3.1) of
\cite{SeppSurvey} and establishes a hydrodynamic limit with a deterministic
concave flux $f_K$; see Theorem 3 there and \cite{SeppHydro}.
Section 4, Problem 4 of \cite{SeppSurvey} asks whether $f_K$ has a flat
segment on an interval about density $K/2$.

Write $H_K$ for the flux of the same process with all rates equal to one.
The homogeneous flux is positive on $(0,K)$, concave, and symmetric about
$K/2$ \cite[Theorem 9.5]{SeppNotes}. Define
\begin{equation}\label{eq:constants}
 M_K=\max_{0\leq\rho\leq K}H_K(\rho)=H_K(K/2),\qquad
 J_K=rM_K,\qquad
 \mu=\frac1r-\EE\!\left[\frac1{\alpha_0}\right].
\end{equation}
Expectation in this formula is with respect to $Q$. The rate bound makes
$\mu$ finite, and a nonconstant law has $\mu>0$.

\begin{theorem}\label{thm:main}
Let $K$ be a positive integer and let $Q$ be a nonconstant probability law
on $[r,1]$ with essential infimum $r>0$. For the totally asymmetric
$K$-exclusion process with independent site rates of law $Q$,
\begin{equation}\label{eq:plateau}
 f_K(\rho)=J_K
 \quad\text{whenever}\quad
 \left|\rho-\frac K2\right|\leq\mu J_K.
\end{equation}
Moreover, $\max_{0\leq\rho\leq K}f_K(\rho)=J_K$. The quantity $\mu J_K$
is a guaranteed positive half-width of a central flat interval.
\end{theorem}

For $K=1$, the homogeneous flux is $H_1(\rho)=\rho(1-\rho)$ and
$M_1=1/4$. In this case \cref{thm:main} is Sly's theorem
\cite[Theorem 1]{Sly}. Sly adds independent exponential delays, conditionally
on the environment, to couple disordered passage times to homogeneous
passage times at rate $r$. Every path to a given column visits the
intervening columns. The mean added delay along those columns creates a
cusp in an upper bound for the passage shape. Flux--shape duality turns
this cusp into a flat interval.

Bahadoran and Bodineau prove a plateau theorem for TASEP with sufficiently
dilute site disorder under their finite-box hypothesis (H), and obtain
quantitative limits for its endpoints and flux
\cite[Theorems 2.4 and 2.6]{BB}. Their discussion in Section 3.2, p.~13 of
arXiv:1602.06718v4, following Theorem 3.4, states that their approach also
applies to other growth models with columnar disorder, including
$K$-exclusion. That paragraph does not state the all-laws finite-$K$
interval \eqref{eq:plateau}.

Let $\tau_\alpha$ denote the passage shape for the disordered process,
and let $\tau_r$ denote its homogeneous rate-$r$ counterpart; their
definitions are given below. For $K\geq2$, the comparison established here is
\begin{equation}\label{eq:intro-comparison}
 \tau_\alpha(x,y)
 \leq \tau_r(x,y)-\mu|x|
 \leq \frac{y+Kx/2}{rM_K}-\mu|x|,
 \qquad y>\max\{0,-Kx\}.
\end{equation}
The finite-capacity growth graph has a vertical edge in addition to the
two edges of the ordinary exclusion representation. This edge preserves
the mandatory-column argument. Reflection symmetry supplies the second
inequality in \eqref{eq:intro-comparison}, without an explicit expression
for $H_K$ or a differentiability assumption on the homogeneous shape.

The proof uses the height and growth framework of \cite{SeppHydro}, its
disordered hydrodynamic formulation in \cite{SeppSurvey}, and the
homogeneous statements and passage construction in \cite{SeppNotes}.
We specify the normalization and inversion used from that framework in
Section~\ref{sec:prelim}. Section~\ref{sec:comparison} proves the
comparison. Section~\ref{sec:proof} determines the maximal current and
deduces \cref{thm:main}. The exact endpoints of the flat segment are not
identified by \eqref{eq:plateau}.

\section{The process and its growth description}\label{sec:prelim}

\subsection{Generator, heights, and flux}
Let $\eta=(\eta_i)_{i\in\ZZ}$ be a configuration in
$\{0,1,\ldots,K\}^{\ZZ}$. When $\eta_i>0$ and $\eta_{i+1}<K$, let
$\eta^{i,i+1}$ be the configuration obtained by decreasing $\eta_i$ by
one and increasing $\eta_{i+1}$ by one. Conditional on $\alpha$, the
generator on local functions $F$ is
\begin{equation}\label{eq:generator}
 L^\alpha F(\eta)=\sum_{i\in\ZZ}\alpha_i
 \mathrm{1}\{\eta_i>0,\ \eta_{i+1}<K\}
 \bigl[F(\eta^{i,i+1})-F(\eta)\bigr].
\end{equation}
The indicator in \eqref{eq:generator} is the only occupation dependence
in the rate. In particular, the rate is not multiplied by either the
number of particles or the number of available places.

Choose heights $h_i(t)$ with
\begin{equation}\label{eq:height}
 h_i(t)-h_{i-1}(t)=\eta_i(t).
\end{equation}
A jump from $i$ to $i+1$ decreases $h_i$ by one and changes no other
height. If $N_i(t)$ is the number of such jumps up to time $t$, then
\begin{equation}\label{eq:current-normalization}
 h_i(t)=h_i(0)-N_i(t).
\end{equation}
Thus the macroscopic equation for a height profile $u$ is
$u_t=-f_K(u_x)$, and its derivative is the conservation law
$\rho_t+\partial_x f_K(\rho)=0$. The empty and full configurations have
zero current, so this normalization gives $f_K(0)=f_K(K)=0$.

For the wedge initial state, take
\begin{equation}\label{eq:wedge}
 \eta_i(0)=K\mathrm{1}\{i\leq0\},\qquad
 h_i(0)=K\min\{i,0\}.
\end{equation}
The hydrodynamic height formula gives a deterministic convex function
$k_\alpha:\RR\to[0,\infty)$ such that, for $t>0$,
\begin{equation}\label{eq:wedge-limit}
 -\frac1n h_{\lfloor nx\rfloor}(nt)
 \longrightarrow G_x(t):=t k_\alpha(x/t)
 \quad\text{in probability}.
\end{equation}
The subscript $\alpha$ indicates the disordered model; the limiting
function depends on $Q$, not on the particular environment. Its dual
relations are
\begin{equation}\label{eq:duality}
 k_\alpha(v)=\sup_{0\leq\rho\leq K}\{f_K(\rho)-v\rho\},\qquad
 f_K(\rho)=\inf_{v\in\RR}\{k_\alpha(v)+v\rho\}.
\end{equation}
These formulas are the sign convention $k_\alpha=-f_K^*$ for the
concave conjugate $f_K^*$ in \cite[equation (3.2)]{SeppSurvey}.
In particular, if
\begin{equation}\label{eq:boundary}
 b(x)=\max\{0,-Kx\},
\end{equation}
then $k_\alpha(v)\geq b(v)$, and
\begin{equation}\label{eq:central-duality}
 k_\alpha(0)=\max_{0\leq\rho\leq K}f_K(\rho).
\end{equation}

The height normalization in \eqref{eq:wedge-limit} can also be read
directly from the density hydrodynamics of
\cite[Theorem 3]{SeppSurvey}. The rightmost occupied site in
\eqref{eq:wedge} advances at rate at most one. For fixed $t$ and $L>t$,
the site $\lfloor nL\rfloor$ therefore remains in the untouched empty
region up to time $nt$ with probability tending to one. On this event,
for $x<L$,
\begin{equation}\label{eq:height-anchor}
 h_{\lfloor nx\rfloor}(nt)
 =-\sum_{i=\lfloor nx\rfloor+1}^{\lfloor nL\rfloor}\eta_i(nt).
\end{equation}
Bounded occupations permit interval indicators in the density limit by
approximation with continuous test functions. Equation
\eqref{eq:height-anchor} fixes the additive constant in the limiting
height. The wedge specialization of the height variational formula is
$t f_K^*(x/t)$, giving \eqref{eq:wedge-limit} and
\eqref{eq:duality} with the integrated-current normalization
\eqref{eq:current-normalization}.

\subsection{The finite-capacity passage graph}
Define the cone
\begin{equation}\label{eq:cone}
 C_K=\{(x,y)\in\RR^2:y\geq0,\ y+Kx\geq0\}.
\end{equation}
On its lattice points use directed steps
\begin{equation}\label{eq:steps}
 (0,1),\qquad (1,0),\qquad (-1,K).
\end{equation}
Conditional on $\alpha$, let $Y_{i,j}$ be independent exponential
random variables of rate $\alpha_i$. For $(i,j)\in C_K\cap\ZZ^2$,
let $\Gamma(i,j)$ be the set of paths from $(0,0)$ to $(i,j)$ using
\eqref{eq:steps}, and set
\begin{equation}\label{eq:passage}
 T_Y(i,j)=\max_{\gamma\in\Gamma(i,j)}
 \sum_{(a,c)\in\gamma\setminus\{(0,0)\}}Y_{a,c}.
\end{equation}
Every such endpoint is reachable. If $i\geq0$, use east and vertical
steps; if $i<0$, use $-i$ northwest steps and then vertical steps.
The grading $2j+Ki$ strictly increases on each edge, so the graph is
acyclic. The bounds below also show that each path set is finite.

To obtain this graph from the process, consider the time at which a
height at site $i$ first reaches an integer level $\ell$. The required
predecessor levels are $\ell+1$ at $i$, $\ell$ at $i-1$, and
$\ell+K$ at $i+1$. Writing $\mathcal T(i,\ell)$ for the completion
time and $W_{i,\ell}$ for the subsequent waiting time gives
\begin{equation}\label{eq:recurrence}
 \mathcal T(i,\ell)=W_{i,\ell}+
 \max\{\mathcal T(i,\ell+1),\mathcal T(i-1,\ell),
                         \mathcal T(i+1,\ell+K)\}.
\end{equation}
The construction of \eqref{eq:recurrence} follows
\cite[Lemma 9.11 and equations (9.40)--(9.44), pp.~148--150]{SeppNotes}.
The conditional independence needed for unequal site rates follows
from the same predecessor ordering, as detailed next.

\begin{lemma}\label{lem:services}
Fix the environment. The waiting times $W_{i,\ell}$ of the eligible
wedge growth cells are mutually independent, and $W_{i,\ell}$ has
exponential rate $\alpha_i$.
\end{lemma}

\begin{proof}
Order the cells so that every predecessor of a cell occurs before it.
Such an ordering exists because the predecessor sets are finite and
the graph is acyclic. Suppose a finite initial segment has already
been considered, and let $(i,\ell)$ be the next cell. Let $S$ be the
maximum of the three predecessor completion times in
\eqref{eq:recurrence}, treating initially completed cells as having
completion time zero. This is a stopping time for the family of
Poisson clocks. At time $S$ the cell becomes eligible. Both neighboring
heights can only decrease, so it remains eligible until the next ring
at site $i$.

Conditional on the fixed environment, the strong Markov property at
$S$ gives independent restarted Poisson clocks with their original
rates, independent of the information up to $S$. The next waiting
time at $i$ is therefore exponential of rate $\alpha_i$, independent
of the past and of all the other restarted clocks. The completion
times in the preceding finite segment are measurable using that
past and those other restarted clocks. Indeed, all cells there at
site $i$ precede $(i,\ell)$ vertically and have already completed
by $S$. No other cell in the segment can depend on the new cell,
because the ordering puts all predecessors first. Its completion
can thus be reconstructed without using any site-$i$ ring after $S$.
Consequently the new waiting time is independent of all preceding
completion times, and hence of their waiting times. Induction proves
the assertion.
\end{proof}

The coordinate change $j=-\ell-1$ moves the first growth cell
$(0,-1)$ to $(0,0)$ and gives precisely \eqref{eq:cone} and
\eqref{eq:steps}. The completion time includes the waiting time at
this first cell. Omitting it in \eqref{eq:passage} removes one
integrable random variable common to every path. The vertical shift
and this omission have no effect on a scaled limit.

For $K=1$, a vertical step can be replaced by an east step followed by
a northwest step. The replacement adds a nonnegative intermediate
weight, so the vertical edge is redundant in the maximizing problem.
For $K>1$, that two-step replacement increases height by $K$, and the
vertical edge must be retained.

\subsection{Shape limits and integrability}
The passage formulation of the height limit is
\begin{equation}\label{eq:shape}
 \frac1n T_Y(\lfloor nx\rfloor,\lfloor ny\rfloor)
 \longrightarrow\tau_\alpha(x,y)
 \quad\text{in probability},\qquad y>b(x).
\end{equation}
All lattice approximations in this paper are taken inside $C_K$.
The deterministic shape is finite, continuous in the interior, and
homogeneous of degree one. The inverse relation at unit time is
\begin{equation}\label{eq:inverse}
 k_\alpha(v)=\inf\{y\geq b(v):\tau_\alpha(v,y)>1\},
\end{equation}
with boundary values interpreted by continuity from the interior.

For completeness, the inversion uses no additional assertion about
stationary measures. The function $G_x$ in \eqref{eq:wedge-limit} is
convex in $t$, since it is the perspective of the convex function
$k_\alpha$. It is nondecreasing, with $G_x(0)=b(x)$, because each
microscopic height decreases. The bound on the site clock rates also
gives $b(x)\leq G_x(t)\leq b(x)+t$, so this extension is continuous
at $t=0$.

The needed comparison with constant rate $r$ is a height comparison.
Use rate-one Poisson clocks with independent uniform marks on $[0,1]$;
the environment-$\alpha$ process accepts a mark at site $i$ when it
is at most $\alpha_i$, while the rate-$r$ process accepts it when it
is at most $r$. Start their heights from the same wedge. The faster
process has lower heights under this coupling. To check preservation
of this order, the only possible violation occurs when the upper
height jumps at a site where the two heights are equal. Eligibility
of the upper height, and the order of both neighboring heights, then
imply eligibility of the lower height. Its clock also accepts the
mark, so it jumps at the same time.

It follows that $G_x(t)\geq G_x^{(r)}(t)$, where the superscript
denotes the homogeneous model. Since $k_r$ is continuous and
$k_r(0)=rM_K>0$, the expression $t k_r(x/t)$ tends to infinity with
$t$. Thus $G_x(t)\to\infty$. Convexity now makes $G_x$
strictly increasing wherever its value exceeds $b(x)$: if
$G_x(s)>b(x)$ and $t>s>0$, then
\[
 G_x(t)\geq G_x(s)+\frac{t-s}{s}\bigl(G_x(s)-b(x)\bigr)>G_x(s).
\]
Thus every level $y>b(x)$ has a unique inverse time. Bracketing this
time by two fixed times in \eqref{eq:wedge-limit} proves
\eqref{eq:shape} and \eqref{eq:inverse}. Continuity of $k_\alpha$
and strictness of the inverse give interior continuity of
$\tau_\alpha$. The identity $G_{cx}(ct)=cG_x(t)$, for $c>0$, gives
its degree-one homogeneity.

We will pass to expectations in \eqref{eq:shape}. The following bound
justifies this step for all the weight fields used below.

\begin{lemma}\label{lem:tails}
Suppose that, conditional on a rate field, the vertex weights are
independent exponentials with rates at least $r$. On every fixed
interior ray of $C_K$, the scaled passage times are uniformly
integrable. Consequently, the limit in \eqref{eq:shape} also holds
in expectation.
\end{lemma}

\begin{proof}
For a path to $(i,j)$, let $a$, $c$, and $e$ count, respectively,
vertical, northwest, and east steps. Then
\[
 j=a+Kc,\qquad i=e-c.
\]
Its length is at most $L=|i|+2j$, and the number of possible paths
is at most $3^{L+1}$. No vertex is repeated. Conditional on the
environment, the exponential moment at parameter $r/2$ of the weight
on one path is at most $2^L$. A union bound gives, for $u\geq0$,
\begin{equation}\label{eq:tail-bound}
 \PP\{T_Y(i,j)>u\}\leq
 3\exp\{L\log6-ru/2\}.
\end{equation}
The bound is uniform in the environment and hence holds after averaging
over it. Along a fixed ray, $L\leq Cn$ for a finite constant $C$.
Putting $u=ns$ in \eqref{eq:tail-bound} gives an integrable exponential
tail uniformly in $n$ once $s>2C\log6/r$. This proves uniform
integrability. Convergence in probability to a constant then implies
convergence in expectation.
\end{proof}

Write $\tau_1$ and $k_1$ for the homogeneous rate-one shapes, and
$\tau_r$ and $k_r$ for the homogeneous rate-$r$ shapes. Scaling the
waiting times gives $\tau_r=\tau_1/r$. At $x=0$, homogeneity and
\eqref{eq:inverse} imply
\begin{equation}\label{eq:reciprocal}
 \tau_1(0,1)=\frac1{k_1(0)}=\frac1{M_K}.
\end{equation}
Here $k_1(0)=M_K$ follows from \eqref{eq:central-duality} and the
homogeneous flux symmetry. Positivity and the upper bound on one site's
clock rate give $0<M_K\leq1$; these properties also appear in
\cite[Theorem 9.5]{SeppNotes}. We use the constant
\begin{equation}\label{eq:A}
 A=\frac{\tau_1(0,1)}r=\frac1{J_K}.
\end{equation}

\section{The comparison argument}\label{sec:comparison}

\subsection{A supporting plane from homogeneous symmetry}
The homogeneous shape is concave on the interior of $C_K$.
One way to see the direction of this inequality is to concatenate two
paths. The corresponding passage values are superadditive, with the
common endpoint counted only once. Translation invariance and the
expectation limits of \cref{lem:tails} give superadditivity of
$\tau_1$. Together with degree-one homogeneity this is concavity.

\begin{lemma}\label{lem:plane}
For every $(x,y)$ in the interior of $C_K$,
\begin{equation}\label{eq:plane}
 \tau_r(x,y)\leq A(y+Kx/2).
\end{equation}
\end{lemma}

\begin{proof}
The linear involution
\[
 S(x,y)=(-x,y+Kx)
\]
preserves the cone and its lattice. It fixes the vertical step and
exchanges $(1,0)$ with $(-1,K)$. Homogeneous independent weights are
invariant in law under $S$, whence
$\tau_1(x,y)=\tau_1(-x,y+Kx)$. Applying concavity at the midpoint of
these two arguments gives
\[
 \tau_1(x,y)\leq\tau_1(0,y+Kx/2)
                 =(y+Kx/2)\tau_1(0,1).
\]
The second coordinate at this midpoint is positive in the interior of
the cone. Divide by $r$ and use \eqref{eq:A}.
\end{proof}

\subsection{Adding exponential delays}
All expectations from now on include both the environment and the
waiting times, unless a conditioning is displayed. Given $\alpha$,
take auxiliary variables independent of the entire field $Y$ and
independent across vertices, with
\begin{equation}\label{eq:auxiliary}
 U_{i,j}=B_{i,j}E_{i,j},\qquad
 B_{i,j}\sim\Ber(1-r/\alpha_i),\qquad
 E_{i,j}\sim\Exp(r).
\end{equation}
Within each pair, $B_{i,j}$ and $E_{i,j}$ are also conditionally
independent. Set $Z_{i,j}=Y_{i,j}+U_{i,j}$.

\begin{lemma}\label{lem:addition}
Conditional on $\alpha$, the field $Z$ consists of independent
rate-$r$ exponential variables, with joint law independent of
$\alpha$. Moreover,
\begin{equation}\label{eq:delay-mean}
 \EE[U_{i,j}\mid\alpha,Y]=d_i,
 \qquad d_i:=\frac1r-\frac1{\alpha_i}\geq0,
 \qquad \EE d_i=\mu.
\end{equation}
\end{lemma}

\begin{proof}
For $s\geq0$, conditional independence gives
\begin{align*}
 \EE[e^{-sZ_{i,j}}\mid\alpha]
 &=\frac{\alpha_i}{\alpha_i+s}
   \left[\frac r{\alpha_i}
    +\left(1-\frac r{\alpha_i}\right)\frac r{r+s}\right]\\
 &=\frac r{r+s}.
\end{align*}
This is the Laplace transform of $\Exp(r)$. The construction gives
conditional independence across vertices, so the conditional joint law
is a product that does not depend on the rates. It remains the same
product law after averaging over the environment. Finally,
\[
 \EE[U_{i,j}\mid\alpha,Y]
 =\left(1-\frac r{\alpha_i}\right)\frac1r=d_i,
\]
which proves \eqref{eq:delay-mean}.
\end{proof}

The fields $Y$ and $U$ share their column environment. Their
independence is conditional on that environment; unconditional
independence is neither asserted nor needed.

\begin{proposition}\label{prop:comparison}
For every interior point $(x,y)$ of $C_K$,
\begin{equation}\label{eq:comparison}
 \tau_\alpha(x,y)\leq\tau_r(x,y)-\mu|x|
                       \leq A(y+Kx/2)-\mu|x|.
\end{equation}
\end{proposition}

\begin{proof}
Fix a lattice endpoint $(i,j)$. Choose a maximizing path
$\gamma_Y\in\Gamma(i,j)$ for $T_Y(i,j)$, resolving any ties by a
fixed deterministic order. This path is measurable with respect to
$\sigma(\alpha,Y)$. With the starting vertex omitted from all sums,
\begin{equation}\label{eq:path-subtract}
 T_Y(i,j)=\sum_{\gamma_Y}Z-\sum_{\gamma_Y}U
           \leq T_Z(i,j)-\sum_{\gamma_Y}U.
\end{equation}
Here $T_Z$ denotes the passage functional \eqref{eq:passage} with
weights $Z$.

Every step changes the column by $-1$, $0$, or $1$. If $i>0$,
every path visits each column in $I_i=\{1,\ldots,i\}$; if $i<0$,
it visits each column in $I_i=\{i,\ldots,-1\}$. Put $I_0=\varnothing$.
Since all $d_a$ are nonnegative, \cref{lem:addition} implies the
pointwise conditional bound
\begin{equation}\label{eq:mandatory-columns}
 \EE\!\left[\sum_{\gamma_Y}U\,\middle|\,\alpha,Y\right]
     =\sum_{(a,c)\in\gamma_Y\setminus\{(0,0)\}}d_a
     \geq\sum_{a\in I_i}d_a.
\end{equation}
The omitted starting vertex causes no loss, since $0\notin I_i$.
Repeated visits to a column only increase the left side. Taking
expectations yields
\begin{equation}\label{eq:finite-comparison}
 \EE T_Y(i,j)\leq\EE T_Z(i,j)-|i|\mu.
\end{equation}
The rightmost sum in \eqref{eq:mandatory-columns} uses fixed columns,
so the path's dependence on the environment introduces no change to
its mean.

Apply \eqref{eq:finite-comparison} to lattice points approaching
$(nx,ny)$ and divide by $n$. By \cref{lem:tails}, both passage terms
converge in expectation to their shapes. By \cref{lem:addition}, the
shape for $Z$ is $\tau_r$. This proves the first inequality in
\eqref{eq:comparison}; \cref{lem:plane} proves the second.
\end{proof}

\section{The maximal current and the flat interval}\label{sec:proof}

\subsection{Long intervals of slow sites}
The comparison above supplies the lower bound on the flux near half
filling. To identify its value, we first determine the passage shape
in the vertical direction.

\begin{lemma}\label{lem:slow-intervals}
Fix an integer $m\geq1$ and $\varepsilon>0$. Almost surely there is an interval
of $2m+1$ consecutive sites on which every rate is at most
$r+\varepsilon$. The interval can be chosen as a measurable function
of $\alpha$ alone.
\end{lemma}

\begin{proof}
The definition of essential infimum gives
$p_\varepsilon:=Q([r,r+\varepsilon])>0$.
Take a deterministic sequence of disjoint intervals, each of length
$2m+1$. Their all-slow events are independent, each of probability
$p_\varepsilon^{2m+1}$. The probability that none of the first $N$
intervals is all slow is
$(1-p_\varepsilon^{2m+1})^N$, which tends to zero. Choosing the first
all-slow interval in this sequence proves both assertions.
\end{proof}

\begin{proposition}\label{prop:capacity}
The vertical passage shape and the maximal flux satisfy
\begin{equation}\label{eq:capacity}
 \tau_\alpha(0,1)=A,\qquad
 k_\alpha(0)=\max_{0\leq\rho\leq K}f_K(\rho)=J_K.
\end{equation}
\end{proposition}

\begin{proof}
For the upper bound on passage times, use independent rate-one
exponential variables $E'_{i,j}$, independent of $\alpha$, and write
$Y_{i,j}=E'_{i,j}/\alpha_i$. Since $\alpha_i\geq r$,
\begin{equation}\label{eq:capacity-upper}
 \tau_\alpha(0,1)\leq\tau_r(0,1)=A.
\end{equation}

For the reverse bound, first fix $m$ and $\varepsilon$. Choose a slow
interval $[a-m,a+m]\cap\ZZ$ as in \cref{lem:slow-intervals}.
Consider all directed paths from $(a,b)$ to $(a,b+m)$, for an integer
$b$ to be chosen. Any such path stays in the selected interval.
Indeed, its northwest-step count is at most $m/K$, and its total
east-step count is equal to its northwest-step count. Its horizontal
displacement at any intermediate vertex therefore has absolute value
at most $m$.

Let $D_m(a,b)$ denote the union of the vertices of these paths.
The waiting times in this set dominate $E'_{i,j}/(r+\varepsilon)$.
Let $V_m$ be the homogeneous rate-one passage value from $(0,0)$ to
$(0,m)$ with its starting weight omitted. On each translated set
$D_m(a,b+\ell m)$, the rate-one field $E'$ therefore supplies a copy
$V_m^{(\ell)}$ of $V_m$.

These copies are independent. To verify this, note first that vertical
coordinates never decrease along a path. Successive sets can intersect
only on their common horizontal row. On its top row, the lower set
has columns at most $a$: any remaining motion to its endpoint must
be eastward. On its bottom row, the upper set has columns at least
$a$: its initial motion on that row is eastward. Thus the only common
vertex is $(a,b+(\ell+1)m)$. It is counted as the terminal vertex of
the lower passage and omitted as the starting vertex of the upper
passage. The weighted vertex sets are disjoint. Conditional on
$\alpha$, translation invariance of $E'$ now gives independent,
identically distributed values $V_m^{(\ell)}$.

The strip may be far from the origin, but its position is fixed at
this stage. Set
\[
 b=K\max\{-a,0\},\qquad d=K\max\{a,0\},\qquad C=b+d.
\]
There is a path from $(0,0)$ to $(a,b)$: use only east steps if
$a\geq0$, and only northwest steps if $a<0$. From $(a,B)$ there is
a path to $(0,B+d)$, by northwest steps if $a>0$ and east steps if
$a\leq0$. These two connectors have fixed height cost $C$.

For $n\geq C$, put $q=\lfloor(n-C)/m\rfloor$. Concatenate the first
connector, $q$ maximizing strip passages, and the return connector,
then fill the remaining height with vertical steps at column zero.
Nonnegative connector weights give
\begin{equation}\label{eq:strip-sum}
 T_Y(0,n)\geq\frac1{r+\varepsilon}
                    \sum_{\ell=0}^{q-1}V_m^{(\ell)}.
\end{equation}
The whole concatenated path is admissible in $C_K$: it is built from
the steps \eqref{eq:steps} starting at the origin. The strip is
chosen using only $\alpha$, so its selection does not bias the field
$E'$.

The conditional strong law and $q/n\to1/m$ show that, almost surely,
\begin{equation}\label{eq:strip-liminf}
 \liminf_{n\to\infty}\frac{T_Y(0,n)}n
       \geq\frac{\EE V_m}{m(r+\varepsilon)}.
\end{equation}
Convergence in probability in \eqref{eq:shape} transfers this bound to
the deterministic value $\tau_\alpha(0,1)$. For example, take an
almost surely convergent subsequence and apply
\eqref{eq:strip-liminf} along it. The homogeneous shape limit and
\cref{lem:tails} give
\[
 \lim_{m\to\infty}\frac{\EE V_m}m=\tau_1(0,1).
\]
Taking integer $m\to\infty$ and then rational
$\varepsilon\downarrow0$ proves
$\tau_\alpha(0,1)\geq\tau_1(0,1)/r=A$.
These are countable limits; the slow-interval and strong-law events
can be intersected without changing their probability.

Together with \eqref{eq:capacity-upper}, this gives the first equality
in \eqref{eq:capacity}. At $x=0$, degree-one homogeneity and
\eqref{eq:inverse} give $k_\alpha(0)=1/\tau_\alpha(0,1)=J_K$.
Equation \eqref{eq:central-duality} gives the remaining assertion.
\end{proof}

The order of limits in this argument is essential to the bound: the
strip width and its position are fixed before $n$ tends to infinity.
Only after that limit is the width increased. No interval with length
proportional to $n$ is required to have all of its rates near $r$.

\subsection{From the passage bound to the flux}
We isolate the geometric implication of the comparison to keep track
of the boundary of the cone.

\begin{lemma}\label{lem:line}
Let $\rho\in[0,K]$ and $J>0$. Suppose that
\[
 \tau_\alpha(v,J-\rho v)\leq1
 \quad\text{whenever }J-\rho v>b(v).
\]
Then $f_K(\rho)\geq J$.
\end{lemma}

\begin{proof}
Fix $v\in\RR$ and put $y_0=J-\rho v$. If $y_0\leq b(v)$, then
$k_\alpha(v)\geq b(v)\geq y_0$. If $y_0>b(v)$, passage times are
nondecreasing in the vertical coordinate, since a path can be extended
by vertical steps with nonnegative weights. The same is true of their
limits. The assumed bound and \eqref{eq:inverse} therefore imply
$k_\alpha(v)\geq y_0$. In both cases,
\[
 k_\alpha(v)+\rho v\geq J.
\]
Take the infimum over $v$ in \eqref{eq:duality}.
\end{proof}

\begin{proof}[Proof of \cref{thm:main}]
First note that the proposed interval is contained in $[0,K]$.
Indeed,
\begin{equation}\label{eq:width-size}
 \mu J_K=\left(1-r\EE\!\left[\frac1{\alpha_0}\right]\right)M_K
             <M_K.
\end{equation}
For $K\geq2$ use $M_K\leq1\leq K/2$; for $K=1$ use $M_1=1/4$.

Take $\rho$ with $|\rho-K/2|\leq\mu J_K=\mu/A$. At every interior
point $(v,J_K-\rho v)$ of $C_K$, \cref{prop:comparison} gives
\begin{align}
 \tau_\alpha(v,J_K-\rho v)
 &\leq A\bigl(J_K+(K/2-\rho)v\bigr)-\mu|v|\notag\\
 &=1+A(K/2-\rho)v-\mu|v|\leq1.\label{eq:line-bound}
\end{align}
The last inequality follows from $A|K/2-\rho|\leq\mu$.
By \cref{lem:line}, $f_K(\rho)\geq J_K$. By
\cref{prop:capacity}, $f_K(\rho)\leq J_K$ at every density. This
proves \eqref{eq:plateau} and the asserted maximum.

Finally, nonconstancy and the essential-infimum condition imply
$\PP\{\alpha_0>r\}>0$. The nonnegative random variable
$1/r-1/\alpha_0$ is positive on this event, so $\mu>0$.
Since $M_K>0$ and $r>0$, the half-width $\mu J_K$ is positive.
\end{proof}

\begin{corollary}\label{cor:binary}
For $Q=p\delta_r+(1-p)\delta_s$ with $0<p<1$ and
$0<r<s\leq1$, where $\delta_a$ denotes unit mass at $a$, the flux
equals $rM_K$ throughout
\[
 \left|\rho-\frac K2\right|
 \leq (1-p)\left(1-\frac rs\right)M_K.
\]
\end{corollary}

\begin{proof}
Substitution of $\EE[1/\alpha_0]=p/r+(1-p)/s$ into
\eqref{eq:width-size} gives the bound.
\end{proof}

\section{Remarks}

The theorem gives a subinterval of the set of maximizing densities.
If
\[
 \rho_-:=\inf\{\rho\in[0,K]:f_K(\rho)=J_K\},\qquad
 \rho_+:=\sup\{\rho\in[0,K]:f_K(\rho)=J_K\},
\]
then concavity makes this set an interval, and the conclusion is
\[
 \rho_-\leq K/2-\mu J_K,
 \qquad \rho_+\geq K/2+\mu J_K.
\]
It does not give equality in these endpoint bounds. Determining the
exact plateau width for general laws, and determining the homogeneous
constant $M_K$ explicitly for $K\geq2$, require information beyond
the supporting-plane comparison.

The positive lower bound on the rates enters both the exponential
addition identity and the passage-time integrability estimate. When
the essential infimum is zero, the quantities $1/r$ and $\mu$ used
here are unavailable. Dependent environments also need separate
analysis: the proof of \cref{lem:slow-intervals} uses independent
blocks, while \cref{prop:comparison} uses deterministic shape limits
and a common mean for the mandatory columns. The present assumptions
ensure both properties without additional conditions.

The state-independent departure clock is part of the result's scope.
If the rate is replaced by $\alpha_i g(\eta_i,\eta_{i+1})$ for a
nonconstant function $g$ on pairs of occupations, the waiting time of
an eligible growth cell can depend on neighboring events. The
independent exponential service field in \eqref{eq:recurrence} then
does not follow from the same construction. Whether central plateaus
persist for such finite-capacity disorder models is a separate
question. Rates attached to particle labels likewise define a
different model from \eqref{eq:generator}.

Finally, a flat hydrodynamic flux does not specify the location or
uniqueness of a microscopic high-density region. The present result
concerns the infinite-volume current--density relation. It does not
determine interface fluctuations or the time required for a finite
system to reach equilibrium.

\section*{Funding}
Computational resources were provided in part by Google Cloud through its
education credits program (credit 530550103).

\section*{Competing interests}
The author declares no competing interests.

\section*{Data availability}
No data were generated or analysed for this paper.

\section*{Declaration of AI use}
This paper is part of a broader research project on AI-assisted proof. The proofs were produced and checked with the
assistance of AI systems (OpenAI GPT-6.1 and Anthropic Claude).

\begin{thebibliography}{9}

\bibitem{BB}
C.~Bahadoran and T.~Bodineau,
\emph{Quantitative estimates for the flux of TASEP with dilute site disorder},
Electron. J. Probab. \textup{23} (2018), paper no.~44, 1--44.
\href{https://doi.org/10.1214/18-EJP137}{doi:10.1214/18-EJP137}.
\href{https://arxiv.org/abs/1602.06718v4}{arXiv:1602.06718v4}.

\bibitem{SeppHydro}
T.~Sepp\"al\"ainen,
\emph{Existence of hydrodynamics for the totally asymmetric simple
$K$-exclusion process},
Ann. Probab. \textup{27} (1999), no.~1, 361--415.
\href{https://doi.org/10.1214/aop/1022677266}{doi:10.1214/aop/1022677266}.

\bibitem{SeppSurvey}
T.~Sepp\"al\"ainen,
\emph{Recent results and open problems on the hydrodynamics of disordered
asymmetric exclusion and zero-range processes},
Resenhas Inst. Mat. Estat. Univ. S\~ao Paulo \textup{4} (1999), no.~1, 1--15.
\href{https://doi.org/10.11606/resimeusp.v4i1.74889}{doi:10.11606/resimeusp.v4i1.74889}.
\href{https://arxiv.org/abs/math/9906194}{arXiv:math/9906194}.

\bibitem{SeppNotes}
T.~Sepp\"al\"ainen,
\emph{Translation invariant exclusion processes},
unpublished book manuscript, December 11, 2008.
\url{https://people.math.wisc.edu/~tseppalainen/excl-book/ajo.pdf}.

\bibitem{Sly}
A.~Sly,
\emph{Note on the flux for TASEP with general disorder},
preprint (2016), \href{https://arxiv.org/abs/1609.06589v1}{arXiv:1609.06589v1}.

\end{thebibliography}
\end{document}
